\section{Why Cologero is Melancholy}

\begin{quotex}
Cornix cornici oculos non effodiet (A crow doesn't rip out the eyes of another crow)

The fullness of existence, life's true richness, does not consist solely in health and happiness but in an ever-expanding range of joy and sorrow; and the broader the range, the richer life becomes.
\flright{\textsc{Valentin Tomberg}}

Some deceive themselves who guard themselves with fear and prudence from every occasion of sin; but when at times they are molested by horrible, filthy, and fearful thoughts, and also sometimes by most loathsome visions, they are confounded and lose heart, and make themselves believe they are forsaken by, and wholly banished from God, and cannot persuade themselves that the Holy Spirit can abide in a mind filled with such thoughts.
\flright{\textsc{Lorenzo Scupoli}, \emph{Spiritual Combat}}
\end{quotex}


I hadn't been in much contact with Cologero until recently, when he asked me to fill in due to his preoccupation with family and health issues. Since he stopped posting translations, he has not needed me so much. Although I have known him for quite some time, our paths have been diverging. He has become a sort of urban sannyasi, isolated and meditative, whereas I plan to retire to Belize where anything goes and the dollar still has clout.

This does not mean that he has not been working on some posts. There are several in different states of development, although I cannot persuade him to bring them to completion. When I proposed this topic to him, he ignored me. Since silence means consent, I am posting it. He did however give me the quotes for the epigraph, which he wanted to be included.

I want to use him to elucidate the influences of \textbf{Saturn}, \textbf{Jupiter}, and \textbf{Venus}, who are still active, even in this age. This is based on the many conversations we've had over the years, especially when he was more Jovial (under the influence of Jupiter) and Venereal (under the influence of Venus).

\paragraph{Soul Sister}

I need to start with something that belongs at the end: the notion of the \emph{âme soeur}. Gornahoor's first post was about the \emph{âme soeur}: Lucinde (in two parts)\footnote{\url{https://gornahoor.net/?p=4}}, and I suspect Cologero's life has been the Platonic quest to reunite with his soul sister. The typical English translation of ``soulmate'' is hardly adequate, since it is just a romanticized name for a lover.

The \emph{âme soeur} (literally ``soul sister'') is not just another life partner, and certainly not a concubine. The idea of ``sister'' should indicate that it is a chaste, not sexual, relationship. Moreover, it indicates a close bond of affinity, not in the biological or genetic sense, but certainly in terms of spiritual races.

There is no doubt that he believes in this literally. He points out the ideal of Courtly love in the Middle ages. This teaching is esoteric as it was the basis of the \emph{Fedeli d'Amore}. Dante and Beatrice are the prime example. Although non-physical and non-matrimonial, even unrequited, the image of Beatrice in his soul was still able to guide \textbf{Dante} through the planets, the Empyrean, to the gates of Heaven. \textbf{Boris Mouravieff} describes courtly love in this way:

\begin{quotex}
Courtly Love is the raison d'être for the couple of and the Lady of his Dreams; without it, their polarity remains spiritualty sterile and they fall back into the common condition. ... In the Middle Ages, the Knight and his Lady, who considered themselves spiritually ONE --- in our terminology, polar beings --- did not venture into marriage. On the contrary, they parted, accepting the risk of never meeting again and knowing that if they did not triumph over a hard test, their love would degenerate, losing its meaning and its marvelous power. They knew that, by separating from each other for an exploit, they stood the chance
\end{quotex}


I was stunned when he told me that he had lost contact with his \emph{âme soeur} around the end of May 2011. When pressed, he claimed that he did not think it odd or even a problem. Then he went off into a discussion of Quantum Entanglements\footnote{\url{https://en.wikipedia.org/wiki/Quantum\_entanglement}}, in which related, but separated, entities still react to each other, across vast distances, faster than the speed of light, instantaneously.

Since the speed of thought is instantaneous, the entanglement of human souls across space and time should not be a surprise. By the Law of Affinity, related souls are brought into the world together for them to play out their destiny.

\paragraph{Saturn}


Of the four humours (Sanguine, Choleric, Phlegmatic, and Melancholy), melancholy is associated with Saturn. This must not be confused with clinical depression. The latter is a medical and psychological condition, whereas the former is metaphysical. Cologero is far from depressed; by profane psychological standards, he has no such issues.

I've taken most of the following quotes from \textbf{Frances Yates}' study on \emph{The Occult Philosophy in the Elizabethan Age}. One of the works she mentions is \textbf{Aristotle}'s \emph{Problemata physica} which discusses melancholy as the humour of heroes and great men.

Saturn is the outermost or highest planet in the cosmic order, and so is nearest to the divine source of being and closest to the angelic hierarchy. Therefore, Saturn is associated with the loftiest contemplations. In the Renaissance, the notion of the melancholy here whose genius is akin to madness became common.

Hence, Saturnians are inspired students and contemplators of highest truths. To be melancholy was a sign of genius; the `gifts' of Saturn, the numbering and measuring studies attributed to the melancholic, were to be cultivated as the highest kind of learning which brought man nearest to the divine.

\textbf{Heroic Frenzy}


Heroic frenzy, or madness, or furor, which according to Plato is the source of all inspiration, when combined with the black bile of the melancholy temperament produces great men; it is the temperament of genius. All outstanding men have been melancholics, heroes like Hercules, philosophers like Empedocles or Plato, and practically all the poets.

The humor melancholicus, when it takes fire and glows, generates the frenzy (furor) which leads us to wisdom and revelation, especially when it is combined with a heavenly influence, above all with that of Saturn ... Therefore Aristotle says in the Problemata that through melancholy some men have become divine beings, foretelling the future like Sybils ... while others have become poets ... and he says further that all men who have been distinguished in any branch of knowledge have generally been melancholics.

\textbf{Angels and Demons}


The following is from \emph{Saturn and Melancholy} by Klibansky, Saxl and Panofsky.

\begin{quotationx}
Moreover, this humor melancholicus has such power that they say it attracts certain demons into our bodies, through whose presence and activity men fall into ecstacies and pronounce many wonderful things . . . This occurs in three different forms corresponding to the threefold capacity of our soul, namely

\begin{itemize}


\item the imagination (\emph{imaginatio})

\item the rational (\emph{ratio})

\item and the mental (\emph{mens})
\end{itemize}

For when set free by the humor melancholicus, the soul is fully concentrated in the imagination, and it immediately becomes a habitation for the lower demons, from whom it often receives wonderful instruction in the manual arts; thus, we see a quite unskilled man suddenly become a painter or an architect, or a quite outstanding master in another art of some kind.

If the demons of this species reveal the future to us, they show us matters related to natural catastrophes and disasters, for instance approaching storms, earthquakes, cloud-bursts, or threats of plague famine and devastation.

But when the soul is fully concentrated in the reason, it becomes the home of the middle demons; thereby it attains knowledge of natural and human things; thus, we see a man suddenly become a philosopher, a physician, or an orator; and of future events they show us what concerns the overthrow of kingdoms and the return of epochs, prophesying in the same way that the Sybil prophesied to the Romans.

But when the soul soars completely to the intellect it becomes the home of the higher demons, from whom it learns the secrets of divine matters, as for instance the law of God, the angelic hierarchy, and that which pertains to the knowledge of eternal things and the soul's salvation; of future events they show us for instance approaching prodigies, wonders, a prophet to come, or the emergence of a new religion, just as the Sybil prophesied Jesus Christ long before lie appeared.
\end{quotationx}


\textbf{Summary}


We can dispense with the idea of the planets controlling the humours and setting them at birth, since we are free beings. However, we can still put ourselves under the influences of the planets at various times. Saturn is the planet of the Sahasrara chakra\footnote{\url{https://gornahoor.net/?p=8618}}, as indicated by \textbf{Johann Gichtel}. Agrippa wrote on the three degrees of melancholy in \emph{De occulta philosophia}.

After I looked into this topic more, I became a little less worried about Cologero. Nevertheless, there is more to the story as the next two sections will show.

\paragraph{Jupiter}


The danger of melancholy needs to be tempered. Frances Yates points out this advice from the Renaissance Platonist, Marsilio Ficino:

\begin{quotex}
Addressing students who were thought to suffer from melancholy through solitariness and concentration on their studies, he advises that the Saturnian or melancholic man should not avoid the deep study to which he is prone by temperament but \emph{should take care to temper the Saturnian severity with Jovial and Venereal influences}.
\end{quotex}


A man under the influence of Jupiter is more sanguine. During those jovial moments, Cologero was fond of quoting Rene Guenon:

\begin{quotex}
There is therefore no occasion for despondency; and even were there no hope of achieving any visible result before the collapse of the modern world through a catastrophe, that would still not be a valid reason to refrain from embarking upon a work extending in scope far beyond the present time. Those who may feel tempted to give way to discouragement should remind themselves that nothing accomplished within the order can ever be lost, that confusion, error and darkness can enjoy no more than a specious and purely ephemeral triumph.
\end{quotex}


True knowledge, or gnosis, therefore cannot lead to hopelessness. An example of the latter is \textbf{H. G. Wells}, a very learned man who became hopeless at the end of his life. He predicted the end of human life in \emph{Mind at the End of its Tether}, not unlike the prophets of Artificial Intelligence do today. That is the end result of secularism and atheism.

Hopelessness is a vice and Hope is a virtue.

The antidote to Saturn is Jupiter and Venus. The Jupiter man is sanguine and those under Venus are gifted with grace and loveliness.

\paragraph{Venus}

\begin{quotex}
Il y a une femme dans toutes les affaires ; aussitôt qu'on me fait un rapport, je dis : « Cherchez la femme ! » 
\flright{\textsc{Alexander Dumas}, \emph{The Mohicans of Paris}}
\end{quotex}

In our youth, Cologero and I used to ``chercher la femme'' together. We were not then as cynical as Dumas, who saw a woman behind every problem, but regarded a woman as the very solution. That is why I urge him to get out more, to yield to Venus who is the source of grace and loveliness.

But I noticed the change in him as he ``got religion'', so to speak. Then, after an affair, he felt the need to confess his adultery or fornication, after which he had to endure the same lecture about leading another into his own sin. Then he would make a half-hearted promise to stop it in other to receive absolution. At night, he would say his beads in expiation. Eventually his fear of the priest became stronger than his desire for sex.

I've tried to convince him that is no way for a man to live. When we last spoke on this topic, about two weeks ago, he reminded me of this thought:

\begin{quotex}
Did you never love a young girl? Were you never commanded by the person beloved to do something which you did not wish to do? have you never flattered your little slave? have you never kissed her feet? And yet if any man compelled you to kiss Caesar's feet, you would think it an insult and excessive tyranny. What else then is slavery? Did you never go out by night to some place whither you did not wish to go, did you not expend that you did not wish to expend, did you not utter words with sighs and groans, did you not submit to abuse and to be excluded?
\flright{\textsc{Epictetus}, \emph{Discourses}}
\end{quotex}


He has a point there. How many outdoor art shows have I had to endure, how often have I spent money I didn't have, just for the hope to kiss the feet of a pretty woman at the end of the day (and I suspect Epictetus is using a euphemism). I admit it, I am a slave to the libido. Cologero threw it right back at me, ``Slavery, even willful slavery, is no way for a man to live.''

He explained that true Love must be freely given and accepted. I countered that the risk of ``abuse and exclusion'' would still exist. He scoffed at me, ``If I wanted unconditional love, I would buy a puppy!''

I let him continue uninterrupted:

\begin{quotex}
Haven't you noticed that, once you've had carnal relations with a woman, how easy it is to repeat it, even if you haven't seen her in years? The physical body has its own intelligence and remembers all its former attractions. That is the consequence of promiscuity.
\end{quotex}


Then he made a strange analogy, even if the conclusion is Traditional:

\begin{quotex}
The sexual bond, when ordered properly, is the analog of a higher bond. Erections are not under a man's conscious control. Yet the mere presence, or even thought, of a woman can raise it up, just as the moon raises the tides. That can be experienced as the infamous ``battle of the sexes'', with the two sides vying for power and control. Or else, it is the experience of non-duality, a Platonic return to androgny, in which the two ``become one flesh'' in their passion.

Without that experience, a fundamental aspect of the world will remain unknown to you. Of course, if that is truly the natural order, then raising erections via artificial stimulation is both disordered and unnatural.
\end{quotex}


Cologero requested that, if I had to write on this topic, to mention \emph{The Faeire Queene} by \textbf{Edmund Spenser}. Since I haven't time to delve into it myself, I will rely again on Frances Yates.

The first book of \emph{The Faerie Queene} is about the \textbf{Red Cross Knight} who represents Holiness and is accompanied by the lady, \textbf{Una}. It is a \emph{solar book}, full of solar imagery.

You see, he can't seem to move away from this theme of Solar Knights and Courtly Love. I presume that the three visions of Venus in Book VI are what he is driving at:

\begin{itemize}


\item First of all the \textbf{planet Venus}, star of love, whose gifts of grace and beauty, when rightly received and not turned by evil will to lust, are so attractive.

\item Then comes the \textbf{Neoplatonic Venus}, representing the beauty of Him through whom all things are beautiful, with many references to Plato, particularly Phaedrus and Symposium.

\item And finally there is the \textbf{angelic Venus}, for Venus is at her best when the angelic hierarchy of the Principalities favours her. Then indeed Venus shines forth, full of grace and charm, courteous and gentle. Such Venereans are agreeable to both God and man.
\end{itemize}

Ultimately, to be somewhat trite, a balance of the Saturnine, Jovial, and Venereal impulses is best.


\flrightit{Posted on 2018-04-20 by Aeneas}

\begin{center}* * *\end{center}

\begin{footnotesize}
\begin{sffamily}
\texttt{rob on 2018-04-20 at 07:23 said:}

``Hopelessness is a vice. Hope is a virtue.'' Isn't this simplification? Each of these compliment each other and lead into one another. It is only after the most sublime hope has wilted from its want of realization that the most melancholic hopelessness forms, and vice versa. There is a middle road where neither hope nor hopelessness is entertained, and this, I think is more where Guenon is coming from in that quote; hence: ``even if there were no hope'' one should continue one's work (because both melancholy and joy in their earthly sense are distractions from it). In other words, to ``balance'' saturnial and jovial influences is to free oneself from both, no? At least, I think it is naive to say that hopelessness is any more a vice than hope. In fact, I think the former is maybe closer to the mark of virtue, a la the pessimism of Schopenhauer, or the idea of resignation from the world in the various religions. To lose hope for the world is to quit the world, which involves both attachment to dread and to hope. What is more conducive to self-resignation: suffering or comfort? I think it has to be the former, because as long as one is in ignorance, there is nothing to stir one to seek after something beyond the limits which suffering reveals. It was not the years of comfort that awoke the Buddha into longing, but the sight of old age and death, and it is not Christ's moments of leisure which sting our hearts with guilt, but his crucifixion.

If you say I am misinterpreting, and you really meant hope, not in the worldly sense, but as a kind of steadfast resolution after the divine, then I would contest the notion that this is really hope, since, if you asked a sage whether he had ``hope'' after heaven or things godly, I think he would merely chuckle, since nothing has ``yet to come'' in heaven and therefore there is no need to hope for anything. This would therefore be like asking a wealthy man whether he hoped for riches. I can only think of hope as rooted in a temporal worldview, and for this reason I am inclined to say that any form of hope is really only detrimental to ``progress'' in personal development. Furthermore, hope should not be confused with anything related to prayer (even though I am aware there could be much disagreement about this), because, in my experience, real ``profit'' from prayer does not stem from a kind of hoping but from active yielding to God, which is a kind of liberation from worldly baggage, which does not feel like waiting or wanting, but is an immediate relief and peace, which in my mind is the farthest thing from hope. It could be described as a lack of it, hence why I said that hopelessness seems to be closer to true attainment, and thus I would not necessarily define hopelessness as the opposite of hope. Despair seems to fit better, and immediately it is clear that despair is a kind of hope or much more closely related to it than freedom from hope which we might now call ``hopelessness'', in the same way that fearlessness is freedom from fear.

\hfill

\texttt{Saladin on 2018-04-20 at 10:41 said:}

``In our youth, Cologero and I used to ``chercher la femme'' together''

That made my day! Had not laughed so hard in a long while! Great job!

\hfill

\texttt{Cologero on 2018-04-20 at 12:47 said:}

@Rob, Hope is an infused (i.e., God given) virtue, so if you don't experience it, it may not be your fault.

\hfill

\texttt{Cologero on 2018-04-20 at 12:50 said:}

@Aeneas, such a banal ending. When you quote others, you sound so good, but your own thoughts are like a Bazooka bubblegum cartoon. I am in an anti-Aristotle mood today. Forget ``balance'': The road of excess leads to the palace of wisdom.

BTW, you don't even meet the moral standards of a crow.

\hfill

\texttt{Arthur Konrad on 2018-04-21 at 11:15 said:}

Excellent as always, but I do have to note, if gentlemen don't mind, that the word ``profane'' is somewhat distasteful and reeks of Masonry.

\hfill

\texttt{manfred arcane on 2018-04-22 at 05:42 said:}

From Halevi's ``The Way of Kabbalah'':

``Occasionally a man becomes too dependent on his maggid, and while it is necessary for the maggid to act as the man's Tiferet for a period, it is not permissible for him to assume this role permanently. While in some branches of the Tradition attachment to the maggid is part of the devotional method, it is not really desirable beyond a certain point, because then the maggid stands in the Way between the man and God, and this is contrary to the Commandments against other gods and bowing down to images. People make a false god of whatever they are devoted to, be it money, status or even a man, however exalted he may be. While the maggid can be the initial link in the chain of connection between Heaven and Earth, he must step aside and allow the man to make his own. It may require him to break up a relationship in which the aspirant's identification with his teacher is either gradually or suddenly dissolved, depending on the circumstance.

The first method is accomplished by sending him away, and the second by demonstrating to the aspirant that the maggid also is human. This is sometimes achieved by the maggid performing an outrageous action aimed at destroying the attached aspirant's image of his teacher.''

\hfill

\texttt{Cologero on 2018-04-26 at 19:24 said:}

Hunh? Profane just means secular as opposed to spiritual. It is the word used by Guenon.

If the Freemasons use ``Profane'' to refer to the uninitiated, then go complain to them.

\end{sffamily}
\end{footnotesize}

